\section{模型假设}

\subsection{基本假设}

为便于建模与求解，本文作出如下假设：

（1）两台设备的时钟运行速率保持一致，设备开机时刻的差异仅表现为时间轴上的固定平移，不考虑时钟漂移和随机时间抖动。

（2）两种定位数据均来源于同一条真实连续轨迹，机器人的真实平面轨迹在观测区间内连续且具有连续的一阶导数，相邻采样时刻之间的轨迹局部变化平滑，因而可以根据离散采样点重建位置及速度变化。

（3）对于含噪声的定位数据，随机测量误差的均值为零，与机器人的真实运动状态相互独立；两种定位方式的随机测量误差彼此独立。

（4）系统偏差在两组数据间表现为固定平移。

（5）机器人同一时刻只能执行一项任务，同一射击目标至多计为一次有效完成，摄影任务可多次完成。

\section{符号说明}

全文主要符号、含义与单位统一列于表\ref{tab:symbols}。

\begin{table}[htbp]
  \centering
  \caption{主要符号说明}
  \label{tab:symbols}
  \begin{tabular}{>{\centering\arraybackslash}p{0.16\textwidth}p{0.57\textwidth}>{\centering\arraybackslash}p{0.13\textwidth}}
    \toprule[1.5pt]
    符号 & 含义 & 单位 \\
    \midrule[1pt]
    $m\in\{1,2\}$ & 定位方式编号 & --- \\
    $t,t_i^{(m)}$ & 参考时刻及方式 $m$ 的第 $i$ 个设备时间戳 & s \\
    $\mathbf r(t),\mathbf z_i^{(m)}$ & 真实二维位置及方式 $m$ 的二维位置观测 & m \\
    $\Delta t_c$ & 问题一中加到方式2时间戳的校正量，$t_{2,\mathrm{ref}}=t_2+\Delta t_c$ & s \\
    $\Delta t$ & 问题二、三的设备时间偏差，$t_{2,\mathrm{ref}}=t_2-\Delta t$ & s \\
    $\mathbf b(t)$ & 方式2减方式1的二维相对系统偏差；常偏差时简记为 $\mathbf b$ & m \\
    $\mathbf x_k,\mathbf P_k$ & 第 $k$ 个事件的状态向量及其协方差矩阵 & 多单位 \\
    $\mathbf R_m,\mathbf Q(\tau;q)$ & 方式 $m$ 的测量噪声协方差及间隔 $\tau$ 下的过程噪声协方差 & 多单位 \\
    $J(\delta),J_H(\delta)$ & 无噪声位置均方误差目标与 Huber 剖面目标 & m$^2$ \\
    $W,I_b$ & HAC--Wald 统计量及是否启用偏差状态的指示变量 & --- \\
    $d_j(t),\theta_j(t)$ & 机器人相对目标 $j$ 的距离与方向角 & m，$^\circ$ \\
    $v(t),a(t)$ & 机器人线速度模与加速度模 & m/s，m/s$^2$ \\
    $\mathcal C,x_c$ & 连续复核后的任务候选集及候选 $c$ 的0--1选择变量 & --- \\
    $I_c,m_c$ & 候选 $c$ 的准备区间及该区间的归一化最小裕度 & s，--- \\
    $N^*,z^*$ & 最大任务数及固定任务数后的最优最小裕度 & --- \\
    \bottomrule[1.5pt]
  \end{tabular}
\end{table}
